\subsection{Eksperimentinio QPU plėtinio projektavimas}
\label{sec:qpu-pletinio-projektavimas}

QPU plėtinys projektuotas kaip eksperimentinis RISC-V specializuotų instrukcijų rinkinys. Visos QPU instrukcijos naudoja \texttt{custom-1} tipo instrukcijų kodą \texttt{0b0101011}, \texttt{funct3 = 1}, o konkreti operacija atskiriama pagal \texttt{funct7}. Tai leidžia QPU instrukcijas dekoduoti tuo pačiu keliu kaip standartines RISC-V instrukcijas, nepakeičiant bendros instrukcijų vykdymo logikos.

Plėtinio būsena padalinta į dvi dalis. Procesoriuje laikomas tik aktyvaus QPU konteksto identifikatorius, pasiekiamas per \texttt{CSR\_QPU\_CTX} adresu \texttt{0x9da}. Visa reali QPU būsena laikoma atskiroje emuliatoriaus QPU posistemėje: joje kiekvienas kontekstas turi savo kvantinius registrus, o patys registrai susiejami su naudojamos kvantinės simuliacijos bibliotekos duomenų struktūromis.

Konteksto valdymas atskirtas nuo kvantinio registro operacijų dėl operacinės sistemos poreikių. Vartotojo programa neturi tiesiogiai kurti globalios QPU būsenos, priklausančios visai sistemai. Vietoj to Linux branduolys gali sukurti kontekstą konkrečiai užduočiai, įrašyti jo identifikatorių į \texttt{CSR\_QPU\_CTX} ir leisti vartotojo programai vykdyti QPU instrukcijas tame kontekste. Šis principas panašus į FPU arba vektorinių registrų būsenos valdymą, kur aparatinė būsena turi būti susieta su vykdomu procesu.
